% =========================================================
% Classical Geometry — Notes Index
% =========================================================

% ---------------------------------------------------------
% Breadcrumb
% ---------------------------------------------------------
\begin{tcolorbox}[
  colback=gray!6,
  colframe=gray!40,
  arc=2pt,
  left=8pt, right=8pt, top=6pt, bottom=6pt,
  title={\small\textbf{Where You Are in the Journey}},
  fonttitle=\small\bfseries
]
\begin{center}
\small
Logic, Sets, Proof Techniques
$\;\to\;$ \textbf{Classical Geometry}
$\;\to\;$ Analytical Geometry
$\;\to\;$ Non-Euclidean Geometry
$\;\to\;$ $\cdots$
\end{center}

\medskip
\noindent\textbf{How we got here.}
Propositional logic and formal proof gave us the tools to reason
rigorously.
Classical geometry is the original domain in which those tools
were deployed: Euclid's \textit{Elements} is the first
sustained axiomatic system in history, and its influence on the
ideal of mathematical proof persists to the present day.

\medskip
\noindent\textbf{What this chapter builds.}
Euclid's five postulates and the synthetic method; congruence
and similarity of triangles; the Pythagorean theorem and its
converses; circle theorems; constructions with compass and
straightedge; area and the method of exhaustion; and the
logical status of the parallel postulate as a genuinely
independent axiom.

\medskip
\noindent\textbf{Where this leads.}
Replacing the parallel postulate with its negation yields
non-Euclidean geometry.
Translating synthetic statements into coordinate equations
yields analytical geometry.
The independence of the parallel postulate, demonstrated
by models of hyperbolic geometry, is also one of the
earliest instances of the model-theoretic reasoning developed
in Volume~I.
\end{tcolorbox}
\vspace{1em}

% ---------------------------------------------------------
% Structural Role
% ---------------------------------------------------------
\begin{tcolorbox}[
  colback=gray!6,
  colframe=gray!40,
  arc=2pt,
  left=8pt, right=8pt, top=6pt, bottom=6pt,
  title={\small\textbf{Structural Role: Classical Geometry — The Axiomatic Template}},
  fonttitle=\small\bfseries
]
Classical geometry is the historical and logical origin of the
axiomatic method.

Euclid's programme --- state your axioms, define your objects,
derive everything else --- is the template that every other
chapter in this project follows.
Working through the \textit{Elements} therefore serves a dual
purpose: it builds geometric intuition and it makes the
axiomatic method visceral and familiar.

The parallel postulate occupies a special place.
For two millennia, mathematicians attempted to derive it from
the other four.
Its independence, established by consistent non-Euclidean
models, is one of the great conceptual milestones of
mathematics and a direct precursor to the model-theoretic
ideas in Volume~I.
\end{tcolorbox}
\vspace{1em}

% ---------------------------------------------------------
% Structural Roadmap
% ---------------------------------------------------------
\subsection*{Structural Roadmap}

\begin{center}
\textbf{Definitions $\longrightarrow$ Main Theorems
$\longrightarrow$ Consequences and Structural Insight}
\end{center}

The global progression is:

\begin{enumerate}
  \item Euclid's postulates and the synthetic method
  \item Congruence criteria: SAS, ASA, SSS
  \item The Pythagorean theorem and its converse
  \item Similarity and proportionality theorems
  \item Circle theorems: inscribed angles, tangents, power of a point
  \item Compass-and-straightedge constructions and impossibility results
  \item Area via decomposition and the method of exhaustion
  \item Hilbert's axioms: a modern repair of Euclid's foundations
  \item The logical status of the parallel postulate
\end{enumerate}

\begin{remark*}[Primary sources]
Euclid, \textit{Elements} (Heath translation);
Hartshorne, \textit{Geometry: Euclid and Beyond};
Moise, \textit{Elementary Geometry from an Advanced Standpoint}.
\end{remark*}

% ---------------------------------------------------------
% Status
% ---------------------------------------------------------
\vspace{1em}
\begin{tcolorbox}[
  colback=gray!6, colframe=gray!40, arc=2pt,
  left=8pt, right=8pt, top=6pt, bottom=6pt,
  title={\small\textbf{Status: Planned}},
  fonttitle=\small\bfseries
]
\begin{center}\Large\bfseries Coming Soon\end{center}
\vspace{6pt}
\noindent Notes, proofs, and exercises will appear here in a future revision.
\end{tcolorbox}

% Future sections (uncomment as content is written):
% \input{volume-v/geometry/classical-geometry/notes/notes-euclid-postulates}
% \input{volume-v/geometry/classical-geometry/notes/notes-congruence}
% \input{volume-v/geometry/classical-geometry/notes/notes-pythagorean}
% \input{volume-v/geometry/classical-geometry/notes/notes-circles}
% \input{volume-v/geometry/classical-geometry/notes/notes-constructions}
% \input{volume-v/geometry/classical-geometry/notes/notes-area}
% \input{volume-v/geometry/classical-geometry/notes/notes-hilbert-axioms}
% \input{volume-v/geometry/classical-geometry/notes/notes-parallel-postulate}
